\chapter{Protocolo de evaluación}
\label{Protocolo de evaluación}

La plantilla reserva este anexo a los cuestionarios empleados durante la fase de evaluación. En este
trabajo \textbf{no procede}: la evaluación no se basa en la valoración subjetiva de usuarios, sino en
la ejecución automatizada de casos de prueba dentro del simulador, con un criterio de éxito medido
por el propio entorno. En su lugar se documenta aquí el protocolo experimental completo, que cumple
la misma función de permitir que un tercero reproduzca la evaluación.

\section{Condiciones de la evaluación}

Parámetros fijos de toda evaluación: entorno, \textit{embodiment}, fondo, presupuesto de episodio,
semillas y número de tiradas. Se justifica cada uno.

\section{Procedimiento}

Secuencia de pasos de una tirada de evaluación, desde el arranque del proceso que sirve la política
hasta la escritura del fichero de resultados.

\section{Registro de resultados}

Estructura del fichero de resultados y campos que se extraen de él para el análisis.

\begin{table*}[ht]
	\centering
	\caption{Campos registrados por episodio}
	\label{tab:campos-registro}
	\makebox[\textwidth]{
		\begin{tabular}{|l|p{10cm}|}
			\cline{1-2}
			Campo & Contenido \\ \hline
			[PENDIENTE] & \\ \hline
		\end{tabular}
	}
\end{table*}

\section{Análisis estadístico}

Procedimiento de análisis: cálculo de las tasas y sus intervalos de Wilson, emparejamiento de las
tiradas entre políticas, prueba exacta de McNemar sobre los casos discordantes y contraste de Fisher
para el efecto de la disposición.
